% Orca ALab architecture report.
% Source: docs/reference/alab-architecture-upstream-trustjs-roadmap.md
% Build:  cd docs/reference && tectonic alab-architecture-upstream-trustjs-roadmap.tex
%
% The body is generated from the Markdown source with Pandoc, then styled here
% so the checked-in TeX remains a standalone, reproducible report artifact.
% Options for packages loaded elsewhere
\PassOptionsToPackage{unicode}{hyperref}
\PassOptionsToPackage{hyphens}{url}
\PassOptionsToPackage{dvipsnames,svgnames,x11names}{xcolor}
\documentclass[
  10pt,
  letterpaper,
]{article}
\usepackage{xcolor}
\usepackage{graphicx}
\usepackage{titlesec}
\usepackage{enumitem}
\usepackage{tabularx}
\usepackage{fancyhdr}
\usepackage{lastpage}
\usepackage{framed}
\usepackage{colortbl}
\usepackage{ragged2e}
\usepackage{needspace}
\definecolor{ink}{HTML}{1B2426}
\definecolor{accent}{HTML}{0C8F86}
\definecolor{muted}{HTML}{5C6A6C}
\definecolor{codebg}{HTML}{F4F1EA}
\definecolor{codekw}{HTML}{0C6E67}
\definecolor{softline}{HTML}{D8E4E2}
\definecolor{shadecolor}{HTML}{F4F1EA}
\usepackage[top=0.78in,bottom=0.78in,left=0.78in,right=0.78in]{geometry}
\usepackage{amsmath,amssymb}
\setcounter{secnumdepth}{-\maxdimen} % remove section numbering
\usepackage{iftex}
\ifPDFTeX
  \usepackage[T1]{fontenc}
  \usepackage[utf8]{inputenc}
  \usepackage{textcomp} % provide euro and other symbols
\else % if luatex or xetex
  \usepackage{unicode-math} % this also loads fontspec
  \defaultfontfeatures{Scale=MatchLowercase}
  \defaultfontfeatures[\rmfamily]{Ligatures=TeX,Scale=1}
\fi
\usepackage{lmodern}
\ifPDFTeX\else
  % xetex/luatex font selection
  \setmainfont[]{Charter}
  \setsansfont[]{Avenir Next}
  \setmonofont[]{Menlo}
\fi
\color{ink}
\titleformat{\section}
  {\sffamily\Large\bfseries\color{ink}}
  {}{0pt}{}
  [\vspace{0.22em}\color{accent}\titlerule]
\titleformat{\subsection}
  {\sffamily\large\bfseries\color{ink}}
  {}{0pt}{}
\titleformat{\subsubsection}
  {\sffamily\normalsize\bfseries\color{codekw}}
  {}{0pt}{}
\titlespacing*{\section}{0pt}{1.9em}{0.65em}
\titlespacing*{\subsection}{0pt}{1.35em}{0.45em}
\titlespacing*{\subsubsection}{0pt}{1.05em}{0.3em}
\setlist[itemize]{leftmargin=1.35em,itemsep=0.18em,topsep=0.35em}
\setlist[enumerate]{leftmargin=1.5em,itemsep=0.22em,topsep=0.35em}
\setlength{\headheight}{15pt}
\fancyhf{}
\fancyhead[L]{\sffamily\footnotesize\color{muted}ORCA ALAB EDITION}
\fancyhead[R]{\sffamily\footnotesize\color{muted}ARCHITECTURE \& TRUSTJS ROADMAP}
\fancyfoot[C]{\sffamily\footnotesize\color{muted}Page \thepage\ of \pageref*{LastPage}}
\renewcommand{\headrulewidth}{0.35pt}
\renewcommand{\headrule}{\hbox to\headwidth{\color{softline}\leaders\hrule height \headrulewidth\hfill}}
\renewcommand{\footrulewidth}{0pt}
\raggedbottom
% Use upquote if available, for straight quotes in verbatim environments
\IfFileExists{upquote.sty}{\usepackage{upquote}}{}
\IfFileExists{microtype.sty}{% use microtype if available
  \usepackage[]{microtype}
  \UseMicrotypeSet[protrusion]{basicmath} % disable protrusion for tt fonts
}{}
\usepackage{setspace}
\makeatletter
\@ifundefined{KOMAClassName}{% if non-KOMA class
  \IfFileExists{parskip.sty}{%
    \usepackage{parskip}
  }{% else
    \setlength{\parindent}{0pt}
    \setlength{\parskip}{6pt plus 2pt minus 1pt}}
}{% if KOMA class
  \KOMAoptions{parskip=half}}
\makeatother
\usepackage{color}
\usepackage{fancyvrb}
\newcommand{\VerbBar}{|}
\newcommand{\VERB}{\Verb[commandchars=\\\{\}]}
\DefineVerbatimEnvironment{Highlighting}{Verbatim}{
  commandchars=\\\{\},fontsize=\footnotesize}
\newenvironment{Shaded}{\begin{snugshade}}{\end{snugshade}}
\newcommand{\AlertTok}[1]{\textcolor[rgb]{1.00,0.00,0.00}{\textbf{#1}}}
\newcommand{\AnnotationTok}[1]{\textcolor[rgb]{0.38,0.63,0.69}{\textbf{\textit{#1}}}}
\newcommand{\AttributeTok}[1]{\textcolor[rgb]{0.49,0.56,0.16}{#1}}
\newcommand{\BaseNTok}[1]{\textcolor[rgb]{0.25,0.63,0.44}{#1}}
\newcommand{\BuiltInTok}[1]{\textcolor[rgb]{0.00,0.50,0.00}{#1}}
\newcommand{\CharTok}[1]{\textcolor[rgb]{0.25,0.44,0.63}{#1}}
\newcommand{\CommentTok}[1]{\textcolor[rgb]{0.38,0.63,0.69}{\textit{#1}}}
\newcommand{\CommentVarTok}[1]{\textcolor[rgb]{0.38,0.63,0.69}{\textbf{\textit{#1}}}}
\newcommand{\ConstantTok}[1]{\textcolor[rgb]{0.53,0.00,0.00}{#1}}
\newcommand{\ControlFlowTok}[1]{\textcolor[rgb]{0.00,0.44,0.13}{\textbf{#1}}}
\newcommand{\DataTypeTok}[1]{\textcolor[rgb]{0.56,0.13,0.00}{#1}}
\newcommand{\DecValTok}[1]{\textcolor[rgb]{0.25,0.63,0.44}{#1}}
\newcommand{\DocumentationTok}[1]{\textcolor[rgb]{0.73,0.13,0.13}{\textit{#1}}}
\newcommand{\ErrorTok}[1]{\textcolor[rgb]{1.00,0.00,0.00}{\textbf{#1}}}
\newcommand{\ExtensionTok}[1]{#1}
\newcommand{\FloatTok}[1]{\textcolor[rgb]{0.25,0.63,0.44}{#1}}
\newcommand{\FunctionTok}[1]{\textcolor[rgb]{0.02,0.16,0.49}{#1}}
\newcommand{\ImportTok}[1]{\textcolor[rgb]{0.00,0.50,0.00}{\textbf{#1}}}
\newcommand{\InformationTok}[1]{\textcolor[rgb]{0.38,0.63,0.69}{\textbf{\textit{#1}}}}
\newcommand{\KeywordTok}[1]{\textcolor[rgb]{0.00,0.44,0.13}{\textbf{#1}}}
\newcommand{\NormalTok}[1]{#1}
\newcommand{\OperatorTok}[1]{\textcolor[rgb]{0.40,0.40,0.40}{#1}}
\newcommand{\OtherTok}[1]{\textcolor[rgb]{0.00,0.44,0.13}{#1}}
\newcommand{\PreprocessorTok}[1]{\textcolor[rgb]{0.74,0.48,0.00}{#1}}
\newcommand{\RegionMarkerTok}[1]{#1}
\newcommand{\SpecialCharTok}[1]{\textcolor[rgb]{0.25,0.44,0.63}{#1}}
\newcommand{\SpecialStringTok}[1]{\textcolor[rgb]{0.73,0.40,0.53}{#1}}
\newcommand{\StringTok}[1]{\textcolor[rgb]{0.25,0.44,0.63}{#1}}
\newcommand{\VariableTok}[1]{\textcolor[rgb]{0.10,0.09,0.49}{#1}}
\newcommand{\VerbatimStringTok}[1]{\textcolor[rgb]{0.25,0.44,0.63}{#1}}
\newcommand{\WarningTok}[1]{\textcolor[rgb]{0.38,0.63,0.69}{\textbf{\textit{#1}}}}
\usepackage{longtable,booktabs,array}
\newcounter{none} % for unnumbered tables
\usepackage{calc} % for calculating minipage widths
% Correct order of tables after \paragraph or \subparagraph
\usepackage{etoolbox}
\makeatletter
\patchcmd\longtable{\par}{\if@noskipsec\mbox{}\fi\par}{}{}
\makeatother
% Allow footnotes in longtable head/foot
\IfFileExists{footnotehyper.sty}{\usepackage{footnotehyper}}{\usepackage{footnote}}
\makesavenoteenv{longtable}
\AtBeginEnvironment{longtable}{\small\renewcommand{\arraystretch}{1.16}}
\setlength{\LTpre}{0.55em}
\setlength{\LTpost}{0.65em}
\arrayrulecolor{softline}
\setlength{\emergencystretch}{3em} % prevent overfull lines
\providecommand{\tightlist}{%
  \setlength{\itemsep}{0pt}\setlength{\parskip}{0pt}}
\usepackage{bookmark}
\IfFileExists{xurl.sty}{\usepackage{xurl}}{} % add URL line breaks if available
\urlstyle{same}
% fallback for those not using the hyperref driver hyperxmp:
\makeatletter
\@ifundefined{xmpquote}{\newcommand{\xmpquote}[1]{#1}}{}
\makeatother
\hypersetup{
  pdftitle={Orca ALab Edition: Architecture, Upstream Comparison, and TrustJS Roadmap},
  pdfauthor={ALab Architecture Review},
  pdfsubject={As-built execution model and strategic migration options},
  pdfkeywords={\xmpquote{Orca, ALab, Electron, TypeScript, Rust,
TrustJS}},
  colorlinks=true,
  linkcolor={accent},
  filecolor={accent},
  citecolor={accent},
  urlcolor={accent},
  pdfcreator={Pandoc and Tectonic}}

\begin{document}
\hypersetup{pageanchor=false}
\begin{titlepage}
\thispagestyle{empty}
\vspace*{0.28in}
\begin{center}
  \includegraphics[width=1.02in]{../../resources/icon-dev.png}\\[1.15em]
  {\sffamily\bfseries\fontsize{28}{32}\selectfont\color{ink}
    Orca ALab Edition\par}
  \vspace{0.45em}
  {\sffamily\fontsize{16}{20}\selectfont\color{accent}
    Architecture, Upstream Comparison,\\
    and TrustJS Roadmap\par}
  \vspace{0.8em}
  {\sffamily\large\color{muted}
    As-built execution model and strategic migration options\par}
  \vspace{1.35em}
  {\color{accent}\rule{0.76\textwidth}{1.3pt}}
\end{center}

\vspace{1.2em}
\begin{tabularx}{\textwidth}{@{}
  >{\sffamily\bfseries\color{muted}}p{0.18\textwidth}
  >{\RaggedRight\arraybackslash}X@{}}
Status & Current-state architecture report and decision roadmap \\[0.45em]
Date & July 21, 2026 \\[0.45em]
ALab snapshot & \texttt{d3ba1bcb17324714a838b2970e04fcd96989e84f}
  \quad \texttt{1.4.147-fork.1} \\[0.45em]
Upstream snapshot & \texttt{ac909c8d837621b5b64c70a18af05fed7eb10fad}
  \quad \texttt{1.4.150-rc.0} \\[0.45em]
Trust snapshot & \texttt{10606c9efa9be25cb798c350a55619234a7ce064} \\[0.45em]
Evidence basis & Pinned source and build-entry analysis; no fresh cross-platform
  application run is claimed. \\
\end{tabularx}

\vfill
\setlength{\fboxsep}{12pt}
\noindent\colorbox{accent!10}{%
  \parbox{\dimexpr\textwidth-2\fboxsep\relax}{%
    \sffamily\small\color{ink}
    \textbf{Preferred direction.} Keep TypeScript/JavaScript as the business-logic
    authoring surface, adopt TrustJS through measured admission gates, and retain
    manual Rust porting for terminal, PTY, crypto, protocol, privileged-boundary,
    and proven performance work.}}

\vfill
\begin{center}
  {\sffamily\small\color{muted}ALab Architecture Review}
\end{center}
\end{titlepage}

\pagenumbering{roman}
\hypersetup{pageanchor=true}
\pagestyle{fancy}
{
\setcounter{tocdepth}{3}
\tableofcontents
}
\clearpage
\pagenumbering{arabic}
\setstretch{1.04}
\section{Executive conclusion}\label{executive-conclusion}

Orca ALab is an Electron application with a React/TypeScript product
shell and several deliberately placed Rust execution islands. It is not
currently a native Rust application. TypeScript still starts Electron,
creates windows, owns the preload and IPC surfaces, coordinates
repositories and sessions, deploys the SSH relay, integrates providers,
and renders the product UI. Rust owns the local terminal daemon, the
aterm terminal engines, a broad but bounded set of pure/business
operations, and selected Git, crypto, persistence, and protocol paths.

The central difference from upstream is the terminal and portable-core
data plane:

\begin{itemize}
\tightlist
\item
  \textbf{ALab:} native Rust daemon + native Node-API add-on + aterm
  WASM + shared Rust dispatch exposed through Node-API and WASM.
\item
  \textbf{Upstream:} forked Node daemon + \texttt{node-pty} +
  xterm.js/headless xterm + TypeScript business logic.
\end{itemize}

Both editions still literally boot an Electron main-process JavaScript
bundle, bridge through a sandboxed preload, and mount a React renderer.
ALab has replaced important engines below that shell; it has not
replaced the shell.

For the forward roadmap, \textbf{option 2 should be the default for
business logic: keep authoring features at the TypeScript/JavaScript
level and adopt TrustJS in stages.} Today that means TrustJS-assisted
differential validation while Node remains the production engine. It
does not yet mean executing Orca under TrustJS. Manual Rust porting
should continue only for the places where Rust is intrinsically
valuable: PTYs, terminal parsing/rendering, crypto, protocol and byte
processing, privileged host boundaries, stable state machines, and
measured performance bottlenecks.

That yields the intended hybrid:

\begin{Shaded}
\begin{Highlighting}[]
\NormalTok{React/Electron UI and feature composition                 TypeScript}
\NormalTok{Pure, provider{-}neutral business logic                     TS/JS → TrustJS track}
\NormalTok{Host effects, provider SDKs, filesystem, SSH adaptation   TypeScript ports/adapters}
\NormalTok{Terminal, PTY, crypto, hot parsers, native persistence    Rust}
\end{Highlighting}
\end{Shaded}

TrustJS is promising but not a production replacement today. Its current
product-level description is an in-development Rust JavaScript
interpreter and differential apparatus. It does not parse TypeScript,
load application modules, provide Node APIs, expose a stable
Orca-callable ABI, ship Node-API/WASM bindings, lower to Trust-IR, or
emit validated native code. Those are roadmap gates, not present-tense
capabilities.

\section{Scope and method}\label{scope-and-method}

This report distinguishes three things that older migration documents
sometimes blend together:

\begin{enumerate}
\def\labelenumi{\arabic{enumi}.}
\tightlist
\item
  \textbf{As-built:} code reached by the current
  product\textquotesingle s build and runtime entry points.
\item
  \textbf{Implemented but not product-wired:} Rust crates, native
  prototypes, or migration artifacts that compile or have tests but are
  not on the normal application path.
\item
  \textbf{Target:} the fully native and Trust-verified architecture
  described in migration plans.
\end{enumerate}

The upstream comparison uses the locally fetched \texttt{upstream/main}
ref. Its reflog records a fetch at 2026-07-21 19:00 PDT. No branch was
checked out and no Internet-only state is assumed. The report is a
source/build trace; it does not claim that a fresh application build or
cross-platform runtime matrix was executed for this review.

\section{1. The as-built ALab
architecture}\label{1-the-as-built-alab-architecture}

\subsection{1.1 Process and trust-boundary
map}\label{11-process-and-trust-boundary-map}

\begin{Shaded}
\begin{Highlighting}[]
\NormalTok{OS launches Electron}
\NormalTok{│}
\NormalTok{├─ Electron main process — compiled TypeScript/JavaScript}
\NormalTok{│  ├─ application lifecycle, windows, IPC, services, providers}
\NormalTok{│  ├─ loads orca\_node.node — Rust Node{-}API dynamic library}
\NormalTok{│  ├─ spawns/connects to orca{-}daemon — standalone Rust process}
\NormalTok{│  ├─ serves runtime RPC to CLI/mobile/web/headless clients}
\NormalTok{│  └─ deploys/manages the TypeScript Node relay on SSH hosts}
\NormalTok{│}
\NormalTok{├─ sandboxed preload — compiled TypeScript/JavaScript}
\NormalTok{│  └─ contextBridge exposes the typed window.api IPC surface}
\NormalTok{│}
\NormalTok{├─ Chromium renderer — React + compiled TypeScript/JavaScript}
\NormalTok{│  ├─ product UI, Monaco, pane lifecycle, feature state}
\NormalTok{│  ├─ Rust Git/dispatch and crypto WASM}
\NormalTok{│  └─ shared render worker}
\NormalTok{│     └─ aterm CPU/GPU WASM → OffscreenCanvas}
\NormalTok{│}
\NormalTok{└─ local terminal daemon — Rust}
\NormalTok{   ├─ authenticated local socket / Windows named pipe}
\NormalTok{   ├─ native PTYs and process lifecycle}
\NormalTok{   ├─ session{-}side aterm headless state and checkpoints}
\NormalTok{   └─ binary/NDJSON event and control protocol}
\end{Highlighting}
\end{Shaded}

The executable entry is still \texttt{out/main/index.js}, declared in
\href{https://github.com/andrewyatesai/orca-alab/blob/d3ba1bcb17324714a838b2970e04fcd96989e84f/package.json\#L11}{\texttt{package.json}}
and built from
\href{https://github.com/andrewyatesai/orca-alab/blob/d3ba1bcb17324714a838b2970e04fcd96989e84f/src/main/index.ts\#L1}{\texttt{src/main/index.ts}}.
The main module binds the Rust dispatch surface before importing
Electron. TypeScript then controls \texttt{app.whenReady()}, constructs
\texttt{BrowserWindow}, and loads either the Vite dev URL or the
packaged renderer HTML. The sandboxed preload publishes
\texttt{electron} and \texttt{api} with \texttt{contextBridge}; the
renderer calls React\textquotesingle s \texttt{createRoot} and renders
\texttt{App}.

\subsection{1.2 What TypeScript owns}\label{12-what-typescript-owns}

TypeScript is the control plane and product layer:

\begin{itemize}
\tightlist
\item
  Electron startup, single-instance behavior, windowing, menus,
  certificates, updates, service initialization, and shutdown;
\item
  preload types and the large IPC API exposed to the sandboxed renderer;
\item
  repository, worktree, session, agent, provider, and feature
  orchestration;
\item
  PTY routing, flow control, batching, environment construction, shell
  policy, recovery, and daemon lifecycle;
\item
  SSH connection management, relay deployment, remote
  filesystem/Git/process effects, and compatibility policy;
\item
  runtime RPC used by the CLI, mobile/web clients, and headless
  \texttt{serve} mode;
\item
  React, pane lifecycle, Monaco, editor/browser/source-control surfaces,
  and user interaction.
\end{itemize}

The \texttt{.ts} and \texttt{.tsx} files are not interpreted directly.
electron-vite/Vite and TypeScript compile or bundle them to JavaScript
for Electron main, preload, Chromium, CLI, web, worker, and relay
contexts. This matters: ALab is a multi-context JavaScript product even
where several algorithms behind those contexts are Rust.

\subsection{1.3 How Rust is delivered and
invoked}\label{13-how-rust-is-delivered-and-invoked}

{\def\LTcaptype{none} % do not increment counter
\begin{longtable}[]{@{}
  >{\raggedright\arraybackslash}p{(\linewidth - 4\tabcolsep) * \real{0.1600}}
  >{\raggedright\arraybackslash}p{(\linewidth - 4\tabcolsep) * \real{0.3600}}
  >{\raggedright\arraybackslash}p{(\linewidth - 4\tabcolsep) * \real{0.4800}}@{}}
\toprule\noalign{}
\rowcolor{accent!12}
\begin{minipage}[b]{\linewidth}\raggedright
Rust form
\end{minipage} & \begin{minipage}[b]{\linewidth}\raggedright
Literal artifact and loader
\end{minipage} & \begin{minipage}[b]{\linewidth}\raggedright
Current role
\end{minipage} \\
\midrule\noalign{}
\endhead
\bottomrule\noalign{}
\endlastfoot
Standalone native process & \texttt{orca-daemon} is built from
\texttt{rust/crates/orca-daemon} and spawned by TypeScript & Local PTYs,
process/session lifetime, session-side terminal state, checkpoints,
socket protocol \\
Node-API dynamic library & \texttt{orca\_node.node} is built from
\texttt{native/orca-node} and loaded into Electron main/CLI & Native
aterm headless API, Git parsers and operations, orchestration store, and
aggregate pure/business dispatch \\
Renderer/worker WASM & aterm CPU and GPU WASM are imported by renderer
code and normally hosted in a shared Web Worker & Per-pane parsing,
terminal state, search/selection, and pixels on
\texttt{OffscreenCanvas} \\
Portable business WASM & \texttt{orca-git-wasm} and crypto WASM are
initialized in sandboxed or relay-safe contexts & Shared dispatch/Git
rules in renderer and SSH relay; browser-safe crypto \\
Direct Rust linking & The daemon links \texttt{orca-pty},
\texttt{orca-terminal}, and \texttt{orca-net} directly, plus
\texttt{orca-winpipe} on Windows & Avoids Node-API on the
daemon\textquotesingle s hot path \\
C ABI/native shell prototype & \texttt{orca-ffi} and
\texttt{native/orca-macos} & Prototype/target scaffold; not the normal
electron-builder product path \\
\end{longtable}
}

The live Rust workspace has 26 members in
\href{https://github.com/andrewyatesai/orca-alab/blob/d3ba1bcb17324714a838b2970e04fcd96989e84f/rust/Cargo.toml\#L12}{\texttt{rust/Cargo.toml}},
but workspace membership is not proof of product execution. The shipped
roots are narrower: the daemon, the aggregate Node-API add-on, aterm
WASM, Git/dispatch WASM, and crypto WASM.

The important business-logic seam is
\href{https://github.com/andrewyatesai/orca-alab/blob/d3ba1bcb17324714a838b2970e04fcd96989e84f/src/shared/orca-dispatch-seam.ts\#L1}{\texttt{orca-dispatch-seam.ts}}.
Main and CLI bind it to Rust through Node-API; renderer and relay
contexts bind the same conceptual registry through WASM. The Rust
registry currently spans roughly 79 pure-domain modules. This lets
TypeScript retain host integration while one Rust implementation
supplies portable rules to several JavaScript runtimes.

The boundary is intentionally coarse. A JavaScript caller sends a module
and a serialized input to Rust and receives a serialized result. This is
a much better FFI shape than crossing Node-API/WASM for every field or
byte, but it also means the serialization contract is part of the
architecture.

\subsection{1.4 Literal development
execution}\label{14-literal-development-execution}

\texttt{pnpm\ dev} executes this chain:

\begin{Shaded}
\begin{Highlighting}[]
\NormalTok{pnpm dev}
\NormalTok{  1. build native/orca{-}node → native/orca{-}node/orca\_node.node}
\NormalTok{  2. build rust/crates/orca{-}daemon → rust/target/release/orca{-}daemon}
\NormalTok{  3. rebuild/check Electron{-}ABI native dependencies}
\NormalTok{  4. prepare isolated dev identity, data directory, CLI wrapper and debug port}
\NormalTok{  5. spawn electron{-}vite dev}
\NormalTok{  6. electron{-}vite builds/watches main, preload, workers and React renderer}
\NormalTok{  7. Electron loads out/main/index.js}
\end{Highlighting}
\end{Shaded}

The chain is visible in
\href{https://github.com/andrewyatesai/orca-alab/blob/d3ba1bcb17324714a838b2970e04fcd96989e84f/package.json\#L42}{\texttt{package.json}}
and
\href{https://github.com/andrewyatesai/orca-alab/blob/d3ba1bcb17324714a838b2970e04fcd96989e84f/config/scripts/run-electron-vite-dev.mjs\#L9}{\texttt{run-electron-vite-dev.mjs}}.
The aterm, Git-dispatch, and crypto WASM bundles are not rebuilt by
every \texttt{pnpm\ dev}; their explicit regeneration commands are
\texttt{build:aterm-wasm}, \texttt{build:relay-wasm}, and
\texttt{build:crypto-wasm}.

\texttt{pnpm\ start} is also not a clean build path. It ensures the
Electron native runtime and starts \texttt{electron-vite\ preview},
assuming the JavaScript bundles and Rust/native artifacts already exist.

electron-vite has separate entries for main-process workers/sidecars,
preload, the main renderer, and the coordinator renderer. The build
therefore preserves Electron\textquotesingle s process boundaries rather
than producing one monolithic JavaScript program.

\subsection{1.5 Literal packaged
execution}\label{15-literal-packaged-execution}

The platform packaging scripts build the TypeScript bundles, native
add-on, Rust daemon, CLI, relay, and web bundle before invoking
electron-builder. The macOS commands additionally build their applicable
native helpers. electron-builder places JavaScript in \texttt{app.asar},
but copies the native add-on and daemon as external resources because
one must be loaded by the dynamic loader and the other must be executed.

At launch:

\begin{Shaded}
\begin{Highlighting}[]
\NormalTok{Orca ALab app executable}
\NormalTok{  → Electron runtime}
\NormalTok{  → app.asar/out/main/index.js}
\NormalTok{  → load Resources/orca\_node.node}
\NormalTok{  → locate/spawn Resources/orca{-}daemon[.exe]}
\NormalTok{  → connect with a random token over a local socket/named pipe}
\NormalTok{  → create sandboxed BrowserWindow + preload}
\NormalTok{  → load out/renderer/index.html}
\NormalTok{  → initialize dispatch/crypto WASM and mount React}
\end{Highlighting}
\end{Shaded}

The product identity is deliberately separate from upstream:
\texttt{com.stablyai.orca.staging} / \texttt{Orca\ ALab\ Edition} in
\href{https://github.com/andrewyatesai/orca-alab/blob/d3ba1bcb17324714a838b2970e04fcd96989e84f/config/electron-builder.config.cjs\#L19}{\texttt{electron-builder.config.cjs}}.

There is one build-consistency gap worth fixing: \texttt{build},
\texttt{build:release}, and \texttt{build:unpack} do not explicitly
rebuild \texttt{orca-daemon}, while the platform packaging commands do.
electron-builder rejects a missing required resource, but a generic
build can consume an existing stale daemon. A single canonical artifact
graph should make the daemon dependency explicit for every packaging
entry point.

\subsection{1.6 Literal local terminal
execution}\label{16-literal-local-terminal-execution}

The normal local control path is:

\begin{Shaded}
\begin{Highlighting}[]
\NormalTok{React pane}
\NormalTok{  → preload IPC: pty:spawn / pty:write / pty:resize}
\NormalTok{  → TypeScript PTY router in Electron main}
\NormalTok{  → DaemonPtyAdapter / DaemonClient}
\NormalTok{  → authenticated daemon control socket}
\NormalTok{  → Rust orca{-}daemon}
\NormalTok{  → Rust orca{-}pty}
\NormalTok{  → native OS PTY}
\NormalTok{  → shell / agent process}
\end{Highlighting}
\end{Shaded}

Output returns on the event path:

\begin{Shaded}
\begin{Highlighting}[]
\NormalTok{native PTY bytes}
\NormalTok{  → Rust daemon read loop}
\NormalTok{  → Rust aterm headless state + checkpoint/session bookkeeping}
\NormalTok{  → daemon event socket (binary frames preferred; NDJSON control)}
\NormalTok{  → TypeScript daemon client + flow control}
\NormalTok{  → Electron IPC pty:data}
\NormalTok{  → preload}
\NormalTok{  → pane transport/controller}
\NormalTok{  → shared render worker}
\NormalTok{  → aterm WASM parses bytes and paints OffscreenCanvas}
\end{Highlighting}
\end{Shaded}

Input and resize travel in the opposite direction. The daemon uses
separate control and event-stream connections with authenticated
hello/version negotiation. The reviewed code\textquotesingle s protocol
is 1020 with minimum compatible version 1018; binary streaming was
introduced at 1020.

The renderer first attempts the shared worker. If that path cannot
initialize, it falls back to in-process aterm GPU and then CPU WASM.
There is no xterm.js rendering fallback in ALab.

The daemon is detached so a normal UI exit can leave local sessions
alive for warm reattachment. If daemon startup fails, there is no
alternate Node daemon; the application degrades to the in-process
TypeScript \texttt{LocalPtyProvider} backed by \texttt{node-pty}. Fresh
terminals still work, but daemon persistence across app exit is lost.

Depending on the active features, one session can be represented by as
many as three independent aterm instances:

\begin{enumerate}
\def\labelenumi{\arabic{enumi}.}
\tightlist
\item
  native aterm in the daemon for session/checkpoint state;
\item
  native aterm through Node-API in Electron main for headless
  query/snapshot paths;
\item
  aterm WASM in the renderer for interaction and pixels.
\end{enumerate}

They share an implementation lineage, not object identity. Feeding
equivalent byte streams and maintaining snapshot/parity tests is
therefore a real correctness, memory, and CPU concern.

\subsection{1.7 SSH is a different
runtime}\label{17-ssh-is-a-different-runtime}

The local Rust daemon does not own SSH PTYs. For remote work, Electron
main uploads a content-hashed CommonJS relay bundle and launches it
under Node on the remote host:

\begin{Shaded}
\begin{Highlighting}[]
\NormalTok{Electron main}
\NormalTok{  → ssh2 / SSH channel}
\NormalTok{  → uploaded relay.js under remote Node}
\NormalTok{  → remote node{-}pty, filesystem, Git and agent processes}
\NormalTok{  → framed channel back to Electron main}
\end{Highlighting}
\end{Shaded}

The relay embeds Rust Git/dispatch WASM for portable parsing and
decisions, but the remote I/O and process host remains TypeScript/Node.
That is an intentional portability property: ALab does not need to build
and upload a native add-on for every SSH host architecture. Any future
business-logic execution policy must preserve this use case; a
desktop-only Node-API solution is incomplete.

\subsection{1.8 CLI and headless
service}\label{18-cli-and-headless-service}

The CLI is also TypeScript compiled to \texttt{out/cli/index.js}.
Packaged wrappers run the Electron binary with
\texttt{ELECTRON\_RUN\_AS\_NODE=1}; the CLI binds the same Rust dispatch
surface where available and talks to the app\textquotesingle s runtime
RPC service.

\texttt{orca\ serve} enters the same Electron main bundle without a
normal main window. It still initializes services, the local PTY
provider/Rust daemon, headless terminal support, and runtime RPC. Some
browser work can still require offscreen Chromium windows. ``Headless''
therefore means no normal UI, not ``no Electron.''

\section{2. How upstream Orca runs}\label{2-how-upstream-orca-runs}

At the reviewed upstream commit, Orca has the same high-level Electron
shape:

\begin{Shaded}
\begin{Highlighting}[]
\NormalTok{OS → Electron main JavaScript}
\NormalTok{   → sandboxed preload}
\NormalTok{   → React renderer}
\end{Highlighting}
\end{Shaded}

Its terminal implementation below that shell is different:

\begin{Shaded}
\begin{Highlighting}[]
\NormalTok{React/xterm.js pane}
\NormalTok{  → preload IPC}
\NormalTok{  → Electron main daemon adapter}
\NormalTok{  → forked daemon{-}entry.js with ELECTRON\_RUN\_AS\_NODE=1}
\NormalTok{  → TypeScript DaemonServer}
\NormalTok{  → node{-}pty}
\NormalTok{  → shell}

\NormalTok{shell output}
\NormalTok{  → node{-}pty}
\NormalTok{  → @xterm/headless state + serialization in Node daemon}
\NormalTok{  → main/preload IPC}
\NormalTok{  → renderer xterm.js + WebGL/DOM/canvas addons}
\end{Highlighting}
\end{Shaded}

Upstream\textquotesingle s \texttt{pnpm\ dev} ensures the Electron
native runtime and starts electron-vite; it does not compile
ALab\textquotesingle s Rust add-on or daemon. Its electron-vite
configuration emits \texttt{daemon-entry.ts} as a separate JavaScript
bundle. Main uses Node\textquotesingle s \texttt{fork()} with the
Electron executable running in Node mode, and the child uses
\texttt{node-pty}. The renderer constructs xterm.js \texttt{Terminal},
Fit, Search, Serialize, Unicode, WebLinks, and WebGL-related addons. The
daemon separately uses \texttt{@xterm/headless} to maintain persistent
terminal state and snapshots.

Upstream has native npm dependencies, including \texttt{node-pty}, but
no tracked \texttt{rust/} product workspace at this ref. Its ordinary
feature and business logic remains TypeScript/JavaScript.

\section{3. ALab versus upstream}\label{3-alab-versus-upstream}

{\def\LTcaptype{none} % do not increment counter
\begin{longtable}[]{@{}
  >{\raggedright\arraybackslash}p{(\linewidth - 4\tabcolsep) * \real{0.1500}}
  >{\raggedright\arraybackslash}p{(\linewidth - 4\tabcolsep) * \real{0.4250}}
  >{\raggedright\arraybackslash}p{(\linewidth - 4\tabcolsep) * \real{0.4250}}@{}}
\toprule\noalign{}
\rowcolor{accent!12}
\begin{minipage}[b]{\linewidth}\raggedright
Dimension
\end{minipage} & \begin{minipage}[b]{\linewidth}\raggedright
ALab edition
\end{minipage} & \begin{minipage}[b]{\linewidth}\raggedright
Upstream Orca
\end{minipage} \\
\midrule\noalign{}
\endhead
\bottomrule\noalign{}
\endlastfoot
Product shell & Electron main + sandboxed preload + React renderer,
mostly TypeScript & Same fundamental shell \\
Local terminal daemon & Standalone Rust \texttt{orca-daemon} & Forked
Node/Electron JavaScript daemon \\
Local PTY & Rust \texttt{orca-pty}; \texttt{node-pty} only for degraded
in-process fallback & \texttt{node-pty} in the Node daemon \\
Session-side terminal & aterm linked into Rust daemon &
\texttt{@xterm/headless} in Node daemon \\
Renderer terminal & aterm CPU/GPU WASM, normally in shared worker +
OffscreenCanvas & xterm.js and addons, including WebGL path \\
Main/CLI native logic & Aggregate Rust Node-API add-on &
TypeScript/JavaScript and native npm dependencies \\
Renderer/relay portable core & Rust WASM dispatch/Git/crypto &
TypeScript/JavaScript \\
SSH host & TypeScript Node relay + \texttt{node-pty}, with Rust WASM for
selected pure logic & TypeScript Node relay + \texttt{node-pty} \\
Failure mode & Rust-daemon failure degrades to in-process TS provider;
no Node-daemon twin & Node daemon is the normal daemon \\
Build complexity & Node/pnpm + Rust + WASM + native add-on +
per-platform resource packaging & Node/pnpm + Electron/native npm
rebuilds \\
Default product compiler for Rust & stable Rust 1.96 product path;
optional Trust toolchain in selected scripts & No first-party Rust
product build at this ref \\
Native desktop shell & SwiftUI/C-ABI prototype exists but is not shipped
& Electron product \\
Verification claim & Terminal and migration-specific evidence; not
whole-app or whole-build Trust verification & Conventional
TypeScript/Rust-free product test surface at this ref \\
\end{longtable}
}

\subsection{3.1 Repository
relationship}\label{31-repository-relationship}

The repositories are still closely related, not independent products:

\begin{itemize}
\tightlist
\item
  merge base: \texttt{d9d939a33b5858495ffb33489a952f1ac9293610};
\item
  ALab-only commits from that base: 883;
\item
  upstream-only commits from that base: 99;
\item
  latest explicit upstream merge in ALab:
  \texttt{b206b314790244e85565a5f1bb54afb11f6e7cdc};
\item
  9,095 paths exist in both snapshots and 7,722 of those have identical
  blobs, approximately 85\% of shared paths.
\end{itemize}

ALab has 9,662 fork-only paths and upstream has 471 upstream-only paths,
but 8,190 of the ALab-only paths are vendored Rust dependencies. Raw
line or path counts therefore wildly overstate how much of the product
was rewritten. The more useful conclusion is architectural: most of the
Electron/product shell is still shared, while ALab has replaced or
wrapped selected lower layers.

This shape has two strategic consequences:

\begin{enumerate}
\def\labelenumi{\arabic{enumi}.}
\tightlist
\item
  continued upstream merging remains valuable because most product work
  is still shared;
\item
  every manual semantic port increases the fork\textquotesingle s merge
  and parity burden, so a module should not move to Rust merely because
  it can.
\end{enumerate}

\subsection{3.2 Current caveats and documentation
drift}\label{32-current-caveats-and-documentation-drift}

\begin{itemize}
\tightlist
\item
  \href{https://github.com/andrewyatesai/orca-alab/blob/d3ba1bcb17324714a838b2970e04fcd96989e84f/docs/rust-migration/architecture.md}{\texttt{docs/rust-migration/architecture.md}}
  is a target architecture, not the current runtime. Sections that
  describe the daemon as TypeScript or the native shell as the product
  are historical or aspirational.
\item
  \texttt{rust/README.md} says 18 workspace crates and records older
  aterm provenance; the live manifest has 26 members and the aterm
  gitlink is \texttt{97b9dcbe5f6cf8619f3228d4367a7dca0ac2ff20}.
\item
  The normal product daemon build explicitly selects stable Rust 1.96
  and clears Trust-only flags. Selected add-on/WASM scripts accept
  \texttt{ORCA\_RUST\_TOOLCHAIN=trust}, but that is not the default
  shipping compiler. Consuming aterm code that has its own proof
  campaign does not make the whole Orca binary Trust-verified.
\item
  Rust-daemon Windows named-pipe code compiles, but its source says real
  Windows end-to-end runtime exercise remains incomplete.
\item
  The daemon release profile uses \texttt{panic=abort}: an uncaught Rust
  panic ends the daemon. Recovery restarts the process, not an individual
  session.
\item
  The Node-API add-on is no longer merely an optional terminal
  acceleration. It carries orchestration persistence and broad dispatch
  logic, making artifact availability effectively a product requirement
  at startup.
\end{itemize}

\section{4. What TrustJS actually is
today}\label{4-what-trustjs-actually-is-today}

Trust is a verification-oriented fork of Rust. TrustJS is a set of Rust
crates inside that repository. Its current execution is:

\begin{Shaded}
\begin{Highlighting}[]
\NormalTok{pinned Test262 JavaScript case}
\NormalTok{├─ Node subprocess + embedded trace driver ─┐}
\NormalTok{├─ Bun subprocess + embedded trace driver ──┼─ compare ObservableTrace}
\NormalTok{├─ trust{-}js{-}sem in{-}process Rust semantics ──┤}
\NormalTok{└─ trust{-}js{-}interp in{-}process Rust engine ──┘}
\end{Highlighting}
\end{Shaded}

The current faithful engine parses ECMAScript Script source into a Rust
AST, walks that AST in a Rust interpreter, and returns either an
\texttt{ObservableTrace} or an explicit \texttt{NoCoverage} refusal. Its
public API is shaped for corpus
cases---\texttt{evaluate\_case(includes,\ body,\ strict)}---rather than
application modules. Each evaluation currently creates and joins a fresh
thread with a 32 MiB stack. The value heap is an arena bounded by
per-evaluation resource caps; there is no long-lived runtime garbage
collector.

The TrustJS crates currently do \textbf{not} depend on Trust-IR, Clean,
\texttt{trustc}, or the autoformalization stack. The faithful
interpreter explicitly resolves and proves nothing. It is valuable
because it refuses unsupported behavior and is differentially measured,
not because it already produces proofs.

\subsection{4.1 Published evidence}\label{41-published-evidence}

The latest published S1c dashboard at the reviewed snapshot reports:

{\def\LTcaptype{none} % do not increment counter
\begin{longtable}[]{@{}
  >{\raggedright\arraybackslash}p{(\linewidth - 2\tabcolsep) * \real{0.7000}}
  >{\raggedleft\arraybackslash}p{(\linewidth - 2\tabcolsep) * \real{0.3000}}@{}}
\toprule\noalign{}
\rowcolor{accent!12}
\begin{minipage}[b]{\linewidth}\raggedright
Evidence
\end{minipage} & \begin{minipage}[b]{\linewidth}\raggedleft
Result
\end{minipage} \\
\midrule\noalign{}
\endhead
\bottomrule\noalign{}
\endlastfoot
Frozen Test262 S0 cases & 35,346 \\
Node/Bun runs & 68,549 \\
Node/Bun trace-equal & 67,510 (98.48\%) \\
Node/Bun divergent runs & 1,037, all classified \\
TrustJS faithful-tier runs covered and equal & 38,406 \\
TrustJS faithful-tier divergent & 0 \\
TrustJS faithful-tier \texttt{NoCoverage} & 30,141 \\
Independent-semantics runs covered and equal & 22,122 \\
Independent-semantics divergent & 0 \\
\end{longtable}
}

The parser\textquotesingle s published gate reports 99.6468\% coverage
on its admitted verdict lane, with disagreements audited and unsupported
inputs explicit. The newest TrustJS source commit adds S1d
regular-expression work, but its commit message says gate verification
is in progress and there is no corresponding published S1d
evidence/ratchet row. S1c is therefore the latest authoritative
execution claim.

These are good engineering results and a useful zero-wrong-trace
discipline. They are not evidence that TrustJS runs 56\% of Orca. S0
intentionally excludes modules, async, shared memory, \texttt{Intl},
\texttt{Temporal}, cross-realm behavior, host GC, and other
application/runtime surfaces. Test262 coverage cannot be projected onto
an Electron application.

\subsection{4.2 Missing product
capabilities}\label{42-missing-product-capabilities}

{\def\LTcaptype{none} % do not increment counter
\begin{longtable}[]{@{}
  >{\raggedright\arraybackslash}p{(\linewidth - 2\tabcolsep) * \real{0.3300}}
  >{\raggedright\arraybackslash}p{(\linewidth - 2\tabcolsep) * \real{0.6700}}@{}}
\toprule\noalign{}
\rowcolor{accent!12}
\begin{minipage}[b]{\linewidth}\raggedright
Needed by Orca
\end{minipage} & \begin{minipage}[b]{\linewidth}\raggedright
TrustJS at this snapshot
\end{minipage} \\
\midrule\noalign{}
\endhead
\bottomrule\noalign{}
\endlastfoot
TypeScript/TSX parsing & Not implemented; current parser accepts
JavaScript \\
Named application modules and imports & M2 scope, not current code \\
Promise/event loop/async/await & M2 scope, not current code \\
\texttt{node:} APIs and npm packages & M2 pure-JS subset planned;
native-addon compatibility explicitly deferred \\
Persistent callable runtime & No product API; current entry is case
evaluation \\
Host callbacks/capability ABI & Not implemented for Orca \\
N-API, WASM, or standalone runtime artifact & No product-consumable
runtime is shipped; the only executable is the corpus-focused
differential harness \\
Long-lived GC/object model & Planned after the current arena/test
model \\
Trust-IR lowering & M3/M4 boundary, not current code \\
Autoformalized native output and install certificates & M4 target, not
current code \\
Cross-platform product CI and packaging & Current harness defaults are
Linux-user-specific; no Orca packaging lane \\
Performance evidence & No TrustJS application benchmarks are
published \\
\end{longtable}
}

The Trust README consequently still classifies TrustJS as design/BUILD
work and not a shipped frontend. That is the correct product-level
description despite the substantial M0/M1 implementation.

\subsection{4.3 TypeScript is not proof
input}\label{43-typescript-is-not-proof-input}

An Orca adoption would need a pinned TypeScript-to-JavaScript transform
before the current TrustJS parser could consume code. The artifact
identity must bind at least:

\begin{itemize}
\tightlist
\item
  TypeScript source hash;
\item
  compiler/transpiler version and exact options;
\item
  emitted JavaScript hash;
\item
  TrustJS revision and runtime build identity;
\item
  admitted capability/effect manifest;
\item
  differential corpus and evidence revision.
\end{itemize}

Type annotations, JSDoc, names, and comments may propose intent; they
cannot assert proof facts or narrow the comparison corpus. This is
Trust\textquotesingle s zero-authority-frontend rule. Differential
equality is bounded behavioral evidence. Only a future Trust-IR/Clean
lane can carry the corresponding proof authority.

\subsection{4.4 Do not conflate TrustJS with earlier port
tooling}\label{44-do-not-conflate-trustjs-with-earlier-port-tooling}

\texttt{trust-ts-embed}, \texttt{trust-formalize}, and
\texttt{tools/ts2rust} are useful precursors, but they are not the
current TrustJS runtime. In particular, ts2rust creates or assesses a
separate Rust candidate, compares it against Node on a finite corpus,
and asks Trust about Rust safety. It is a migration assistant, not
general TypeScript semantics and not an integrated Orca execution
engine.

\section{5. Roadmap options}\label{5-roadmap-options}

\subsection{5.1 Decision matrix}\label{51-decision-matrix}

{\def\LTcaptype{none} % do not increment counter
\begin{longtable}[]{@{}
  >{\raggedright\arraybackslash}p{(\linewidth - 4\tabcolsep) * \real{0.1900}}
  >{\raggedright\arraybackslash}p{(\linewidth - 4\tabcolsep) * \real{0.4050}}
  >{\raggedright\arraybackslash}p{(\linewidth - 4\tabcolsep) * \real{0.4050}}@{}}
\toprule\noalign{}
\rowcolor{accent!12}
\begin{minipage}[b]{\linewidth}\raggedright
Criterion
\end{minipage} & \begin{minipage}[b]{\linewidth}\raggedright
Option 1: continue manual Rust porting
\end{minipage} & \begin{minipage}[b]{\linewidth}\raggedright
Option 2: TrustJS-first hybrid
\end{minipage} \\
\midrule\noalign{}
\endhead
\bottomrule\noalign{}
\endlastfoot
Can ship new product logic today & Yes, after port/parity work & Yes in
TypeScript; TrustJS shadowing only today \\
Feature iteration speed & Lower for every manually translated module &
Highest: one maintained TS/JS source \\
Native performance certainty today & High for well-designed ports &
Unknown until later native tier; interpreter unbenchmarked \\
Existing Orca integration seam & Mature Node-API, daemon, WASM, dispatch
& Must build callable/runtime and packaging seams \\
Browser/renderer/SSH portability & Established through WASM and TS relay
& Requires a WASM or portable artifact before primary remote use \\
Behavioral-fidelity cost & Per-module translated tests and dual-run
corpus & Central engine corpus plus per-export admission corpus \\
Proof status today & Rust is not automatically Trust-verified; proof
lane remains separate & Differential evidence only; native/proof
authority is M4+ \\
Upstream merge burden & Grows with every deleted/replaced TS module &
Lower while feature source remains TS/JS \\
Main risk & Feature capacity consumed by translation and semantic drift
& Betting product cutover on an immature runtime if gates are skipped \\
Best fit & Data plane, privileged boundaries, measured hot paths &
Pure/provider-neutral business logic and fast feature work \\
\end{longtable}
}

\subsection{5.2 Option 1 --- continue manual TypeScript-to-Rust
porting}\label{52-option-1--continue-manual-typescript-to-rust-porting}

This option extends the pattern ALab already proved: define a coarse
boundary, port the behavior and tests, expose it through
Node-API/WASM/daemon protocols, dual-run against the TypeScript
implementation, cut over, then remove the hand-maintained twin after an
explicit stability window.

\subsubsection{Phase R1 --- make current ownership
truthful}\label{phase-r1--make-current-ownership-truthful}

\begin{itemize}
\tightlist
\item
  Generate an as-built inventory from actual import/build roots, not
  crate existence or old migration checklists.
\item
  Classify every candidate as UI, host/effect adapter, pure business
  rule, performance data plane, or unused/test-only Rust port.
\item
  Fix the generic-build daemon dependency and establish macOS, Linux,
  Windows, WSL, and SSH artifact matrices.
\end{itemize}

\textbf{Gate:} every shipped Rust root is reproducibly built, packaged,
loaded, and smoke-tested on its applicable hosts; documentation names
current ownership.

\subsubsection{Phase R2 --- harden the Rust substrate already in
production}\label{phase-r2--harden-the-rust-substrate-already-in-production}

\begin{itemize}
\tightlist
\item
  Exercise Windows named-pipe daemon startup, reconnect, upgrade, crash,
  and uninstall paths on real Windows.
\item
  Test daemon protocol downgrade/upgrade and token/socket isolation.
\item
  Measure the cost of multiple aterm state instances and eliminate
  unnecessary duplicate processing where semantics permit.
\item
  Preserve the Node relay/WASM shape for SSH rather than assuming a
  desktop add-on exists remotely.
\end{itemize}

\textbf{Gate:} crash/reconnect/session-persistence and cross-host parity
pass with bounded resources and actionable telemetry.

\subsubsection{Phase R3 --- port only justified vertical
slices}\label{phase-r3--port-only-justified-vertical-slices}

Require at least one concrete reason per module:

\begin{itemize}
\tightlist
\item
  measured CPU, latency, allocation, or startup improvement;
\item
  privileged memory/process boundary;
\item
  protocol/parser totality or security value;
\item
  stable cross-context implementation needed in native and WASM forms;
\item
  mature state machine whose churn is low enough to recover the port
  cost.
\end{itemize}

Freeze a versioned request/result schema first. Capture TypeScript
behavior with unit, property, fuzz, recorded-production, and mutation
corpora. Run both implementations, including provider, GitLab/GitHub,
SSH, WSL, and Git 2.25 compatibility cases where relevant.

\textbf{Gate:} zero unexplained divergence, acceptable performance, all
applicable host contexts green, and an atomic rollback path before
primary cutover.

\subsubsection{Phase R4 --- decide separately whether to pursue a native
shell}\label{phase-r4--decide-separately-whether-to-pursue-a-native-shell}

The existing SwiftUI/C-ABI code is not a consequence of porting one more
business module. A native shell requires its own product decision,
especially for Monaco/editor capability, accessibility, i18n, Linux, and
Windows. Do not let a backend-port program silently become a full UI
rewrite.

\subsubsection{Option 1 assessment}\label{option-1-assessment}

This is deployable and appropriate for the existing Rust data plane. As
the default for ordinary business features, however, it spends
engineering effort twice---first to implement the feature, then to
translate and prove parity---and it increases the cost of every upstream
merge. It should be a selective tool, not the organizing principle for
the entire application.

\subsection{5.3 Option 2 --- TrustJS-first hybrid
(preferred)}\label{53-option-2--trustjs-first-hybrid-preferred}

The practical policy is \textbf{TrustJS-assisted now, TrustJS-executed
selectively later, and TrustJS-native only after its native-validation
gates exist}. Choosing this option does not remove the Rust daemon,
aterm, crypto, or other successful native subsystems. It changes how new
and changing business logic is authored and admitted.

\subsubsection{Phase T0 --- define a TrustJS-ready business island
now}\label{phase-t0--define-a-trustjs-ready-business-island-now}

Keep production execution in Node/Electron and keep writing features in
TypeScript. Move eligible logic into packages with explicit constraints:

\begin{itemize}
\tightlist
\item
  deterministic, provider-neutral functions or reducers;
\item
  no DOM, Electron, filesystem, Git process, network, clock, random,
  locale, or ambient environment access;
\item
  effects represented as typed inputs, outputs, or commands handled by a
  host adapter;
\item
  bounded execution and explicit error values;
\item
  no dependence on module-level singleton state.
\end{itemize}

Good candidates are policy calculations, normalizers, serializers,
reducers, provider-neutral review/workspace decisions, and view-model
projections. React components, Electron services, SSH transport, Git
execution, PTY handling, and provider SDK adapters are not candidates.

Use the existing coarse \texttt{orca-dispatch} idea as the architectural
precedent, but version the new ABI. Either constrain v1 values to a
validated JSON-safe domain or use a tagged value encoding. Plain JSON
silently loses JavaScript distinctions such as \texttt{undefined},
\texttt{NaN}, \texttt{-0}, BigInt, and lone UTF-16 surrogates.

Pin and hash the TS-to-JS transform. Prefer erasable TypeScript and a
modern JS target; avoid TS constructs whose emitted helper/runtime
semantics obscure the source contract.

\textbf{Gate T0:} selected modules are effect-free by enforced
dependency rules, have a versioned ABI, run unchanged under normal
Node/Electron, and carry source/emitted-JS identity plus cross-platform
tests. No TrustJS claim is made yet.

\subsubsection{Phase T1 --- productize TrustJS as a consumable
engine}\label{phase-t1--productize-trustjs-as-a-consumable-engine}

This is work in Trust, not a path dependency from Orca into an entire
mutable rustc checkout. Produce a small, versioned, pinned runtime
artifact with:

\begin{itemize}
\tightlist
\item
  persistent engine lifetime rather than one 32 MiB thread per call;
\item
  module/function invocation or an equivalent registered-export API;
\item
  tagged input/output values and typed \texttt{NoCoverage}/fault
  responses;
\item
  time, recursion, heap, output, and cancellation limits;
\item
  deterministic engine identity and evidence metadata;
\item
  configurable Node/Bun oracle discovery;
\item
  clean locked builds and continuous macOS/Linux/Windows tests;
\item
  cold-start, throughput, memory, and bundle-size benchmarks;
\item
  a standalone/process-isolated runner first, followed by Node-API and
  WASM bindings only if their measurements justify them.
\end{itemize}

Land the pending S1d evidence before using those semantics. Continue the
M1 ratchet: every admitted behavior is exact or refused, never guessed.

\textbf{Gate T1:} a clean pinned artifact builds reproducibly on all
Orca targets, has published coverage and performance evidence, and
exposes a stable callable ABI. This still makes no proof or sandbox
claim.

\subsubsection{Phase T2 --- shadow selected Orca
exports}\label{phase-t2--shadow-selected-orca-exports}

For each T0 module, run the emitted JavaScript through both the ordinary
Node engine and the TrustJS runner in CI and development replay jobs.
Node remains authoritative. Compare tagged results and effects over:

\begin{itemize}
\tightlist
\item
  existing unit and integration cases;
\item
  generated/property and mutation cases;
\item
  recorded, privacy-reviewed production shapes;
\item
  boundary strings, Unicode, numeric corners, malformed data, and
  resource limits;
\item
  macOS, Linux, Windows, WSL, and SSH-relevant data shapes.
\end{itemize}

\texttt{NoCoverage}, timeout, or engine fault means ``not admitted; stay
on Node.'' A behavioral divergence fails the admission gate and produces
a minimized case; it is never silently classified by the source module
itself.

\textbf{Gate T2:} zero divergence and zero \texttt{NoCoverage} for every
admitted export on its complete pinned corpus, deterministic results on
all target OSes, and negative controls that prove the comparator catches
a wrong result.

\subsubsection{Phase T3 --- selective faithful-tier primary
execution}\label{phase-t3--selective-faithful-tier-primary-execution}

Promote one pure module at a time only after T2. Use the same source and
ABI so rollback is engine selection, not a data migration or
hand-maintained code fork. Retain sampled Node shadow evaluation for a
bounded rollout window.

Promotion must be context-aware:

\begin{itemize}
\tightlist
\item
  Electron main/CLI may use a process-isolated or Node-API artifact;
\item
  renderer and web require WASM;
\item
  an export needed on an SSH host stays on remote Node until a portable
  WASM artifact or explicitly supported remote binary exists;
\item
  no module may assume the desktop architecture when its feature runs
  through the relay.
\end{itemize}

\textbf{Gate T3:} module ABI stability, acceptable measured
latency/memory/startup, packaging and signing on every applicable
platform, crash/resource isolation, two release-equivalent windows of
shadow parity, and an exercised atomic rollback.

\subsubsection{Phase T4 --- effectful runtime adoption after TrustJS
M2}\label{phase-t4--effectful-runtime-adoption-after-trustjs-m2}

Do not move async orchestration merely because pure functions work. Wait
for TrustJS\textquotesingle s actual M2 gate: reactor and Promise
semantics, module graph, the admitted pure-JS \texttt{node:} surface,
capability enforcement, async/module differential ledgers, and
reproducible cross-platform packaging.

Even then, keep Electron DOM/rendering, native add-ons, PTYs,
filesystem, Git, SSH, network, and provider APIs behind explicit host
ports. Consider moving only orchestration whose effects can be injected
and audited.

\textbf{Gate T4:} Trust\textquotesingle s published M2 exit conditions
plus Orca-specific disconnect, cancellation, retry, provider, WSL, SSH,
and resource tests.

\subsubsection{Phase T5 --- autoformalized/native tier after TrustJS
M4}\label{phase-t5--autoformalizednative-tier-after-trustjs-m4}

At the future M4 boundary, TypeScript/JavaScript remains zero-authority
input. TrustJS may produce an inspectable Rust+Lean/Trust-IR artifact
and validated native installation candidate. Replace a manual Rust port
only when the generated artifact clears both independent fidelity
evidence and its stated proof/certificate gate.

The maintained source should remain TypeScript/JavaScript; generated
Rust/Trust-IR should be reproducible build artifacts, not a second
source tree edited by hand. That is the feature-velocity payoff of
option 2.

\section{6. Recommended ownership
policy}\label{6-recommended-ownership-policy}

Adopt option 2 as the strategic default with an explicit option-1
carve-out:

{\def\LTcaptype{none} % do not increment counter
\begin{longtable}[]{@{}
  >{\raggedright\arraybackslash}p{(\linewidth - 2\tabcolsep) * \real{0.3400}}
  >{\raggedright\arraybackslash}p{(\linewidth - 2\tabcolsep) * \real{0.6600}}@{}}
\toprule\noalign{}
\rowcolor{accent!12}
\begin{minipage}[b]{\linewidth}\raggedright
New or changing work
\end{minipage} & \begin{minipage}[b]{\linewidth}\raggedright
Default home
\end{minipage} \\
\midrule\noalign{}
\endhead
\bottomrule\noalign{}
\endlastfoot
React, UI behavior, Electron integration, editor/browser views &
TypeScript \\
Pure business rules, reducers, normalizers, provider-neutral planning &
TypeScript/JavaScript in the TrustJS-ready island \\
Filesystem/network/provider/Git/SSH effects and SDK adaptation &
TypeScript host ports unless a native boundary is independently
justified \\
PTY, terminal engine/rendering, crypto, binary protocols, hot parsers &
Rust \\
Stable, measured, privileged, memory-sensitive state machines & Rust or
future admitted TrustJS-native artifact \\
SSH-host pure rules & TS now; TrustJS only with portable WASM/remote
support \\
\end{longtable}
}

The operating rules are:

\begin{enumerate}
\def\labelenumi{\arabic{enumi}.}
\tightlist
\item
  \textbf{Stop broad manual business-logic porting.} A new Rust port
  needs a measured performance, security, platform, or stability reason.
\item
  \textbf{Do not unwind successful Rust infrastructure.} The daemon,
  aterm, crypto, protocol, and established portable-core work are the
  substrate for either strategy.
\item
  \textbf{Keep shipping features in TypeScript.} Shape eligible code so
  it can be admitted by TrustJS later without blocking the feature now.
\item
  \textbf{Treat \texttt{NoCoverage} as a routing fact.} It means stay on
  Node, not waive the case or infer equivalence.
\item
  \textbf{Keep proof language precise.} Node/TrustJS equality is bounded
  differential evidence. TrustJS source and TS annotations have zero
  proof authority. Native/proof claims wait for the actual
  Trust-IR/Clean lane.
\item
  \textbf{Preserve all execution hosts.} Desktop, renderer, web, CLI,
  WSL, SSH, macOS, Linux, and Windows must be part of module admission.
\item
  \textbf{Keep upstream mergeability visible.} Prefer stable seams and
  retained high-level source over fork-wide semantic rewrites.
\end{enumerate}

\section{7. Immediate next decisions}\label{7-immediate-next-decisions}

\begin{enumerate}
\def\labelenumi{\arabic{enumi}.}
\tightlist
\item
  Ratify the ownership table above as the default rule for new work.
\item
  Produce a generated as-built port inventory: \texttt{wired},
  \texttt{shadow-only}, \texttt{test-only}, \texttt{prototype}, or
  \texttt{target} for every Rust crate/module.
\item
  Fix the generic packaging graph so every package build rebuilds or
  verifies the exact daemon artifact it ships.
\item
  Select three small T0 pilots: one normalizer, one reducer/state
  transition, and one provider-neutral policy calculation with strong
  existing tests.
\item
  Define \texttt{orca-business-abi-v1}, including tagged/safe values,
  errors, resource limits, artifact identity, and capability-free
  declarations.
\item
  In Trust, finish and publish the S1d gate, then design the persistent
  callable runner and cross-platform artifact. Do not link Orca directly
  to \texttt{\textasciitilde{}/trust}.
\item
  Add shadow comparison only after T0/T1 gates; keep Node authoritative
  until each export independently passes T2 and T3.
\end{enumerate}

This sequence preserves current feature velocity, continues to exploit
Rust where it already wins, and turns TrustJS maturity into incremental
optionality rather than a blocker or a leap of faith.

\section{Evidence index}\label{evidence-index}

Key ALab sources:

\begin{itemize}
\tightlist
\item
  build and entry points:\par
  \href{https://github.com/andrewyatesai/orca-alab/blob/d3ba1bcb17324714a838b2970e04fcd96989e84f/package.json}{\texttt{package.json}}\\
  \href{https://github.com/andrewyatesai/orca-alab/blob/d3ba1bcb17324714a838b2970e04fcd96989e84f/electron.vite.config.ts}{\texttt{electron.vite.config.ts}}\\
  \href{https://github.com/andrewyatesai/orca-alab/blob/d3ba1bcb17324714a838b2970e04fcd96989e84f/config/scripts/run-electron-vite-dev.mjs}{\texttt{run-electron-vite-dev.mjs}}\\
  \href{https://github.com/andrewyatesai/orca-alab/blob/d3ba1bcb17324714a838b2970e04fcd96989e84f/config/electron-builder.config.cjs}{\texttt{electron-builder.config.cjs}}
\item
  Electron/preload/renderer:\par
  \href{https://github.com/andrewyatesai/orca-alab/blob/d3ba1bcb17324714a838b2970e04fcd96989e84f/src/main/index.ts}{\texttt{src/main/index.ts}}\\
  \href{https://github.com/andrewyatesai/orca-alab/blob/d3ba1bcb17324714a838b2970e04fcd96989e84f/src/main/window/createMainWindow.ts}{\texttt{createMainWindow.ts}}\\
  \href{https://github.com/andrewyatesai/orca-alab/blob/d3ba1bcb17324714a838b2970e04fcd96989e84f/src/preload/index.ts}{\texttt{src/preload/index.ts}}\\
  \href{https://github.com/andrewyatesai/orca-alab/blob/d3ba1bcb17324714a838b2970e04fcd96989e84f/src/renderer/src/main.tsx}{\texttt{src/renderer/src/main.tsx}}
\item
  daemon:\par
  \href{https://github.com/andrewyatesai/orca-alab/blob/d3ba1bcb17324714a838b2970e04fcd96989e84f/src/main/daemon/daemon-init.ts}{\texttt{daemon-init.ts}}\\
  \href{https://github.com/andrewyatesai/orca-alab/blob/d3ba1bcb17324714a838b2970e04fcd96989e84f/src/main/daemon/daemon-spawner.ts}{\texttt{daemon-spawner.ts}}\\
  \href{https://github.com/andrewyatesai/orca-alab/blob/d3ba1bcb17324714a838b2970e04fcd96989e84f/rust/crates/orca-daemon/src/lib.rs}{\texttt{orca-daemon}}\\
  \href{https://github.com/andrewyatesai/orca-alab/blob/d3ba1bcb17324714a838b2970e04fcd96989e84f/rust/crates/orca-daemon/src/protocol.rs}{\texttt{daemon\ protocol}}
\item
  terminal renderer:\par
  \href{https://github.com/andrewyatesai/orca-alab/blob/d3ba1bcb17324714a838b2970e04fcd96989e84f/src/renderer/src/lib/pane-manager/aterm/aterm-strategy-select.ts}{\texttt{aterm-strategy-select.ts}}\\
  \href{https://github.com/andrewyatesai/orca-alab/blob/d3ba1bcb17324714a838b2970e04fcd96989e84f/src/renderer/src/lib/pane-manager/aterm/aterm-shared-render-worker.ts}{\texttt{aterm-shared-render-worker.ts}}\\
  \href{https://github.com/andrewyatesai/orca-alab/blob/d3ba1bcb17324714a838b2970e04fcd96989e84f/src/renderer/src/lib/pane-manager/aterm/aterm-render-worker.ts}{\texttt{aterm-render-worker.ts}}
\item
  Rust seams:\par
  \href{https://github.com/andrewyatesai/orca-alab/blob/d3ba1bcb17324714a838b2970e04fcd96989e84f/native/orca-node/src/lib.rs}{\texttt{native/orca-node}}\\
  \href{https://github.com/andrewyatesai/orca-alab/blob/d3ba1bcb17324714a838b2970e04fcd96989e84f/rust/crates/orca-dispatch/src/lib.rs}{\texttt{orca-dispatch}}\\
  \href{https://github.com/andrewyatesai/orca-alab/blob/d3ba1bcb17324714a838b2970e04fcd96989e84f/src/shared/orca-dispatch-seam.ts}{\texttt{orca-dispatch-seam.ts}}\\
  \href{https://github.com/andrewyatesai/orca-alab/blob/d3ba1bcb17324714a838b2970e04fcd96989e84f/src/relay/git-wasm.ts}{\texttt{relay\ Git\ WASM}}
\item
  SSH:\par
  \href{https://github.com/andrewyatesai/orca-alab/blob/d3ba1bcb17324714a838b2970e04fcd96989e84f/src/relay/relay.ts}{\texttt{relay.ts}}\\
  \href{https://github.com/andrewyatesai/orca-alab/blob/d3ba1bcb17324714a838b2970e04fcd96989e84f/src/main/ssh/ssh-relay-deploy.ts}{\texttt{ssh-relay-deploy.ts}}.
\end{itemize}

Reproducible upstream inspections at the pinned ref:

\begin{Shaded}
\begin{Highlighting}[]
\FunctionTok{git}\NormalTok{ show upstream/main:package.json}
\FunctionTok{git}\NormalTok{ show upstream/main:electron.vite.config.ts}
\FunctionTok{git}\NormalTok{ show upstream/main:src/main/daemon/daemon{-}init.ts}
\FunctionTok{git}\NormalTok{ show upstream/main:src/main/daemon/daemon{-}entry.ts}
\FunctionTok{git}\NormalTok{ show upstream/main:src/main/daemon/pty{-}subprocess.ts}
\FunctionTok{git}\NormalTok{ show upstream/main:src/main/daemon/headless{-}emulator.ts}
\FunctionTok{git}\NormalTok{ show upstream/main:src/renderer/src/lib/pane{-}manager/pane{-}dom{-}creation.ts}
\FunctionTok{git}\NormalTok{ show upstream/main:src/renderer/src/lib/pane{-}manager/pane{-}lifecycle.ts}
\end{Highlighting}
\end{Shaded}

Key Trust evidence at the pinned snapshot:

\begin{itemize}
\tightlist
\item
  \texttt{\textasciitilde{}/trust/docs/design/2026-07-20-trust-native-javascript-engine.md};
\item
  \texttt{\textasciitilde{}/trust/docs/design/2026-07-20-trustjs-m0-scope.md};
\item
  \texttt{\textasciitilde{}/trust/docs/design/2026-07-21-trustjs-m1-scope.md};
\item
  \texttt{\textasciitilde{}/trust/docs/design/2026-07-22-trustjs-m2-scope.md};
\item
  \texttt{\textasciitilde{}/trust/crates/trust-js-interp/src/lib.rs};
\item
  \texttt{\textasciitilde{}/trust/crates/trust-js-differential/src/heads.rs};
\item
  \texttt{\textasciitilde{}/trust/tests/js262/dashboard.md};
\item
  \texttt{\textasciitilde{}/trust/tests/js262/coverage.toml}.
\end{itemize}

\end{document}
